\documentclass[11pt]{article}
\usepackage[paperwidth=220mm,paperheight=320mm,left=16mm,right=16mm,top=18mm,bottom=18mm,headheight=16pt,headsep=10pt,footskip=14pt]{geometry}
\usepackage{fix-cm}
\usepackage{fontspec,amsmath,amssymb,graphicx,array,multirow,tikz}
\usetikzlibrary{decorations.pathreplacing}
\usepackage[unicode,hidelinks]{hyperref}
\usepackage{polyglossia}
\setmainlanguage{latin}
\setotherlanguages{arabic,greek,hebrew,syriac}
\setmainfont{Linux Libertine O}[Ligatures=TeX]
\setsansfont{Linux Biolinum O}[Ligatures=TeX]
\newfontfamily\greekfont{FreeSerif}[Script=Greek]
\newfontfamily\arabicfont{Amiri}[Script=Arabic]
\newfontfamily\hebrewfont{Frank Ruhl Hofshi}[Script=Hebrew]
\newfontfamily\syriacfont{Segoe UI Historic}[Script=Syriac]
\newfontfamily\CapFont{Noto Serif}
\newfontfamily\MoonFont{FreeSerif}
\newcommand{\MoonSym}{\text{\MoonFont ☾}}\newcommand{\SunSym}{\text{\MoonFont ⊙}}
\newfontfamily\EthiopicFont{Ebrima}
\newcommand{\textethiopic}[1]{{\EthiopicFont #1}}
\newfontfamily\NallinoSigns{NallinoSigns.otf}[Path=./fonts/]
\newcommand{\AbjadZero}{{\NallinoSigns\symbol{"E000}}}
\hypersetup{pdftitle={Al-Battani, Opus astronomicum. Nallino, versio Latina et adnotationes (Pars I). S04},pdfauthor={Al-Battani; Carlo Alfonso Nallino (Latin translation and notes)}}
\makeatletter
\def\ps@sourceedition{%
\def\@oddhead{\vbox{\hbox to\textwidth{\fontsize{8}{10}\selectfont C. A. NALLINO — VERSIO ET ADNOTATIONES\hfil\leftmark}\vskip3pt\hrule height.25pt}}%
\let\@evenhead\@oddhead
\def\@oddfoot{\hbox to\textwidth{\fontsize{7}{9}\selectfont\rightmark\hfil\thepage}}%
\let\@evenfoot\@oddfoot}
\makeatother
\pagestyle{sourceedition}
\setlength{\parindent}{0pt}\setlength{\parskip}{0pt}
\newcommand{\BeginSource}[2]{\clearpage\markboth{PAG. #2}{#1}\hypertarget{#1}{}\pdfbookmark[0]{#1 — p. #2}{bk-#1}\fontsize{11.1}{13.8}\selectfont}
\tracinglostchars=2
\newcommand{\Name}[1]{{\addfontfeatures{LetterSpace=4.0}#1}}
\newcommand{\Indent}{\hspace*{1.6em}}
% an Arabic run inside a Latin line: neither its descenders nor its vowel signs above (shadda, damma) open the line
% (the printed lines are evenly spaced)
\newcommand{\ArabicRun}[1]{\smash{\textarabic{#1}}}
\newcommand{\qfrac}[2]{{}^{#1}\!/_{{#2}}}
\newcommand{\sa}{\textsuperscript{\textit{a}}}\newcommand{\sm}{\textsuperscript{\textit{m}}}\newcommand{\sd}{\textsuperscript{\textit{d}}}
\newcommand{\RawBlock}[1]{\par\vspace{3pt}\noindent#1\par\vspace{3pt}}
\newcommand{\CenterLine}[1]{\makebox[\linewidth][c]{#1}}
\newlength{\FullSourceWidth}\newlength{\GutterOffset}\newif\ifNumbersLeft
\newcommand{\DSource}[2]{\BeginSource{AB01-PDF#1}{#2}%
 \setlength{\FullSourceWidth}{\textwidth}\setlength{\GutterOffset}{0pt}%
 \ifodd#2\relax\NumbersLefttrue\else\NumbersLeftfalse\fi}
\newsavebox{\DLBox}
\newcommand{\DLine}[4]{%
 \par\noindent\hypertarget{#1}{}%
 \ifx\relax#2\relax\else
   \ifNumbersLeft
    \rlap{\kern-\GutterOffset\llap{\fontsize{7}{8}\selectfont #2\hspace{3mm}}}%
   \else
    \rlap{\kern\dimexpr\FullSourceWidth-\GutterOffset\relax\hspace{3mm}{\fontsize{7}{8}\selectfont #2}}%
   \fi
 \fi
 \ifx\relax#3\relax\else
   \ifNumbersLeft
    \rlap{\kern\dimexpr\FullSourceWidth-\GutterOffset\relax\hspace{3mm}{\fontsize{7}{8}\selectfont #3}}%
   \else
    \rlap{\kern-\GutterOffset\llap{\fontsize{7}{8}\selectfont #3\hspace{3mm}}}%
   \fi
 \fi
 \sbox{\DLBox}{\strut #4}%
 \ifdim\wd\DLBox>\linewidth\typeout{DIPLOMATICWIDTH|#1|\the\wd\DLBox|\the\linewidth}\fi
 \usebox{\DLBox}\strut\par}
\begin{document}
\thispagestyle{empty}
\begin{center}\Large AL-BATTĀNĪ\par\vspace{4mm}\LARGE OPUS ASTRONOMICUM\par
\vspace{5mm}\large Caroli Alphonsi Nallino versio Latina et adnotationes (pars I)\par\vspace{8mm}
\Large S04 --- editio diplomatica\par\vspace{4mm}\large Paginae impressae 219--219\end{center}
\vspace{12mm}\noindent Singulae lineae paginarum impressarum singulis lineis huius editionis respondent, cum
divisionibus vocabulorum, signis et generibus litterarum (rectis, inclinatis, distantibus). Numeri linearum et
paginae textus Arabici in marginibus stant ut impressi sunt. Textus ex imaginibus paginarum transcriptus est;
litterae Graecae et Arabicae, numeri et signa in imaginibus amplificatis lecta sunt.
Haec transcriptio a nullo homine recognita est.\par
\DSource{0308}{219}
\par\vspace*{\dimexpr 0.08mm-\baselineskip\relax}\prevdepth=0pt
\DLine{AB01-PDF0308-body-L001}{}{}{\mbox{}\rlap{\hspace*{0.00mm}\smash{\rlap{Iidem, Damasci, anno 201 Iezdegerdis, 217 heg. (832 Chr.), teste \Name{Ibn Yūnus} l. c.}}}\rlap{\hspace*{134.36mm}\smash{\llap{82}\rlap{°}}}\rlap{\hspace*{140.29mm}\smash{\llap{1}\rlap{′}}}\rlap{\hspace*{145.47mm}\smash{\llap{37}\rlap{″}}}}
\DLine{AB01-PDF0308-body-L002}{}{}{Tres fratres \Name{Banū Mūsà ibn Shākir} supputabant apogeum primo Iezdegerdis}
\DLine{AB01-PDF0308-body-L003}{}{}{\mbox{}\rlap{\hspace*{7.61mm}\smash{\rlap{anno (inc. 16 Iun. 632) fuisse, ut \Name{Ibn Yūnus} cap. V, \Name{Caussin} 150, narrat}}}\rlap{\hspace*{134.36mm}\smash{\llap{80}\rlap{°}}}\rlap{\hspace*{140.29mm}\smash{\llap{44}\rlap{′}}}\rlap{\hspace*{145.47mm}\smash{\llap{19}\rlap{″}}}}
\DLine{AB01-PDF0308-body-L004}{}{}{\Name{Aḥmad}, unus eorum trium, anno 220 Iezdegerdis (inc. 23 Apr. 851 Chr.), teste}
\DLine{AB01-PDF0308-body-L005}{}{}{\mbox{}\rlap{\hspace*{7.61mm}\smash{\rlap{\Name{Ibn Yūnus}, l. c., et cap. VIII (ap. \Name{Delambre}, p. 99)}}}\rlap{\hspace*{134.36mm}\smash{\llap{84}\rlap{°}}}\rlap{\hspace*{140.29mm}\smash{\llap{33}\rlap{′}}}}
\DLine{AB01-PDF0308-body-L006}{}{}{\mbox{}\rlap{\hspace*{0.00mm}\smash{\rlap{Familia \Name{Banū Amāǵūr} (\Name{Ibn Yūnus} cap. V, \Name{Caussin} p. 152)}}}\rlap{\hspace*{134.36mm}\smash{\llap{83}\rlap{°}}}\rlap{\hspace*{140.29mm}\smash{\llap{35}\rlap{′}}}}
\DLine{AB01-PDF0308-body-L007}{}{}{\mbox{}\rlap{\hspace*{0.00mm}\smash{\rlap{\Name{Mufliḥ ibn Yūsuf}, libertus illius familiae (\Name{Ibn Yūnus} l. c.)}}}\rlap{\hspace*{134.36mm}\smash{\llap{84}\rlap{°}}}\rlap{\hspace*{140.29mm}\smash{\llap{5}\rlap{′}}}}
\par\vspace{1.14mm}
\hypertarget{AB01-PDF0308-P01}{}
\DLine{AB01-PDF0308-body-L008}{}{}{\Indent Addantur duae illae \Name{Persicae} observationes (?) quas, ex opere Ibn Yūnus sumptas,}
\DLine{AB01-PDF0308-body-L009}{}{}{retuli p. 209, adn. 2; posterior, anni 630 Chr. circiter, 80° praebuit ut tabulae \textit{ash–Shāh} (e}
\DLine{AB01-PDF0308-body-L010}{}{}{Pehlevico seu Persico aevi medii in Arabicum translatae) et prior \textit{Sūrya–siddhānta.}}
\par\vspace{3.40mm}
\hypertarget{AB01-PDF0308-H11}{}
\DLine{AB01-PDF0308-body-L011}{}{}{\CenterLine{\textbf{\textit{d}}) — \textit{Pag. 46, l. 7.}}}
\par\vspace{3.36mm}
\hypertarget{AB01-PDF0308-P02}{}
\DLine{AB01-PDF0308-body-L012}{}{}{\Indent Longitudinem lineae EF servavi ut est in cod. Escurialensi et in versione Platonis, nempe}
\DLine{AB01-PDF0308-body-L013}{}{}{61\textsuperscript{\textit{p}} 48′ 35″. Exactius erat:}
\par\vspace{0.35mm}
\hypertarget{AB01-PDF0308-H14}{}
\DLine{AB01-PDF0308-body-L014}{}{}{\CenterLine{$EF = \sqrt{EM^2 + FM^2} = \sqrt{(61^{p}\,48'\,2'')^2 + (1^{p}\,2'\,22''\,\qfrac{1}{2})^2} =$}}
\hypertarget{AB01-PDF0308-H15}{}
\DLine{AB01-PDF0308-body-L015}{}{}{\CenterLine{$= \sqrt{(222482'')^2 + (3472'',5)^2} = \sqrt{49\;498\;240\;324^{\text{IV}},00 + 14\;006\;306^{\text{IV}},25} =$}}
\par\vspace{-0.14mm}
\hypertarget{AB01-PDF0308-H16}{}
\DLine{AB01-PDF0308-body-L016}{}{}{\CenterLine{$= \sqrt{49\;512\;246\;630^{\text{IV}},25} = 222\;513'',47... = 61^{p}\,48'\,33'',47.$}}
\par\vspace{3.61mm}
\hypertarget{AB01-PDF0308-H17}{}
\DLine{AB01-PDF0308-body-L017}{}{}{\CenterLine{\textbf{\textit{e}}) — \textit{Pag. 46, l. 25–26.}}}
\par\vspace{2.19mm}
\hypertarget{AB01-PDF0308-P03}{}
\DLine{AB01-PDF0308-body-L018}{}{}{\Indent Longitudo lineae EG ut in cod. et apud Platonem, nempe 58\textsuperscript{\textit{p}} 12′ 34″; sed p. 47, l. 32,}
\par\vspace{0.14mm}
\DLine{AB01-PDF0308-body-L019}{}{}{cod. rectius secundas 32″ habet. Nam}
\par\vspace{0.21mm}
\hypertarget{AB01-PDF0308-H20}{}
\DLine{AB01-PDF0308-body-L020}{}{}{\CenterLine{$EG = \sqrt{KG^2 + EK^2} = \sqrt{(1^{p}\,2'\,22''\,\qfrac{1}{2})^2 + (58^{p}\,11'\,58'')^2} =$}}
\par\vspace{-0.12mm}
\hypertarget{AB01-PDF0308-H21}{}
\DLine{AB01-PDF0308-body-L021}{}{}{\CenterLine{$= \sqrt{(3742'',5)^2 + (209\;518'')^2} = \sqrt{14\;006\;306^{\text{IV}},25 + 43\;897\;792\;324^{\text{IV}}} =$}}
\hypertarget{AB01-PDF0308-H22}{}
\DLine{AB01-PDF0308-body-L022}{}{}{\CenterLine{$= \sqrt{43\;911\;798\;625^{\text{IV}},25} = 209\;551'',4... = 58^{p}\,12'\,31'',4.$}}
\par\vspace{3.51mm}
\hypertarget{AB01-PDF0308-H23}{}
\DLine{AB01-PDF0308-body-L023}{}{}{\CenterLine{\textbf{\textit{f}}) — \textit{Pag. 48, l. 2-16.}}}
\par\vspace{-\ht\strutbox}\noindent\begin{minipage}[t]{0.391\FullSourceWidth}\vspace{0pt}\hypertarget{AB01-PDF0308-body-F01}{}% Nallino, Part I p. 219: the figure of the equation of the centre (after Almagest III, 5), redrawn from the print:
% the circle ABC about E, the epicycles about A and about B; the straight lines HC (through M, A and E), MF, FA,
% FE, QE (through B, K and P), BG, GK and GE, meeting where the measured strokes meet; the diameter NT through E
% in dashes, each where it printed; the bold italic letters. Measured on the scan (_intake/latin/tools/fig219).
\begin{tikzpicture}[x=0.9928pt,y=-0.9928pt,line cap=butt,every node/.style={font=\fontsize{11}{11}\bfseries\itshape\selectfont,inner sep=0pt}]
\useasboundingbox (0,0) rectangle (212.00,210.00);
\draw[line width=0.83pt] (127.48,121.79) circle[radius=73.67];  % the circle ABC about E
\draw[line width=0.84pt] (127.39,47.68) circle[radius=32.77];  % the epicycle about A
\draw[line width=0.90pt] (54.69,112.59) circle[radius=32.89];  % the epicycle about B
\draw[line width=0.77pt] (127.40,14.92) -- (127.57,195.46);  % HC through M, A and E
\draw[line width=0.74pt] (127.40,18.98) -- (142.70,18.71);  % MF
\draw[line width=0.64pt] (142.70,18.71) -- (127.42,48.09);  % FA
\draw[line width=0.71pt] (142.70,18.71) -- (127.65,121.91);  % FE
\draw[line width=0.84pt] (22.02,108.79) -- (127.65,121.91);  % QE through B, K and P
\draw[line width=0.73pt] (54.53,112.70) -- (82.47,94.98);  % BG
\draw[line width=0.70pt] (82.47,94.98) -- (79.30,115.90);  % GK
\draw[line width=0.75pt] (82.47,94.98) -- (127.65,121.91);  % GE
% the dashes of the diameter NT
\draw[line width=0.91pt,line cap=round] (78.33,173.45)--(79.91,171.87) (81.19,170.67)--(82.31,169.34) (83.63,168.13)--(85.11,166.45) (86.62,164.82)--(87.91,163.43) (89.47,161.85)--(90.30,160.97) (91.79,159.46)--(93.02,158.13) (94.41,156.67)--(95.74,155.33) (97.01,153.96)--(98.25,152.65);
\draw[line width=0.91pt,line cap=round] (99.51,151.41)--(100.73,150.08) (102.14,148.54)--(103.50,147.11) (104.86,145.79)--(106.58,143.99) (107.89,142.62)--(109.36,141.05) (110.72,139.64)--(112.01,138.25) (113.31,136.93)--(114.89,135.23) (116.19,133.93)--(117.77,132.23) (119.00,130.93)--(120.42,129.44);
\draw[line width=0.91pt,line cap=round] (121.73,128.17)--(122.97,126.85) (124.11,125.62)--(125.26,124.57) (129.42,119.91)--(130.56,118.73) (131.82,117.46)--(133.44,115.92) (134.45,114.82)--(135.22,114.00) (136.46,112.71)--(137.55,111.49) (138.64,110.38)--(139.96,109.02) (141.19,107.79)--(142.37,106.39);
\draw[line width=0.91pt,line cap=round] (143.59,105.27)--(144.77,103.87) (145.94,102.70)--(147.30,101.27) (148.52,99.99)--(149.71,98.74) (150.98,97.46)--(152.12,96.29) (153.38,94.90)--(154.74,93.47) (156.14,92.02)--(157.50,90.59) (158.80,89.22)--(160.41,87.55) (161.79,86.12)--(163.15,84.68);
\draw[line width=0.91pt,line cap=round] (164.32,83.46)--(165.93,81.79) (167.10,80.65)--(168.24,79.33) (169.51,77.96)--(171.11,76.40) (172.50,75.01)--(174.03,73.39) (175.03,72.34)--(176.44,70.82);
\node at (126.84,8.24) {\normalfont\fontsize{11}{11}\bfseries\itshape\selectfont \resizebox{7.73pt}{7.14pt}{H}};
\node at (144.38,12.26) {\normalfont\fontsize{11}{11}\bfseries\itshape\selectfont \resizebox{5.83pt}{7.49pt}{F}};
\node at (120.93,23.20) {\normalfont\fontsize{11}{11}\bfseries\itshape\selectfont \resizebox{8.38pt}{7.14pt}{M}};
\node at (122.39,53.54) {\normalfont\fontsize{11}{11}\bfseries\itshape\selectfont \resizebox{6.95pt}{7.48pt}{A}};
\node at (185.13,65.48) {\normalfont\fontsize{11}{11}\bfseries\itshape\selectfont \resizebox{5.07pt}{7.61pt}{T}};
\node at (64.38,75.14) {\normalfont\fontsize{11}{11}\bfseries\itshape\selectfont \resizebox{3.45pt}{7.01pt}{I}};
\node at (84.95,88.90) {\normalfont\fontsize{11}{11}\bfseries\itshape\selectfont \resizebox{6.70pt}{8.21pt}{G}};
\node at (15.88,109.10) {\normalfont\fontsize{11}{11}\bfseries\itshape\selectfont \resizebox{6.21pt}{7.42pt}{Q}};
\node at (49.45,117.80) {\normalfont\fontsize{11}{11}\bfseries\itshape\selectfont \resizebox{6.27pt}{7.60pt}{B}};
\node at (76.83,121.48) {\normalfont\fontsize{11}{11}\bfseries\itshape\selectfont \resizebox{7.62pt}{7.48pt}{K}};
\node at (91.07,122.96) {\normalfont\fontsize{11}{11}\bfseries\itshape\selectfont \resizebox{5.64pt}{7.38pt}{P}};
\node at (133.08,124.76) {\normalfont\fontsize{11}{11}\bfseries\itshape\selectfont \resizebox{6.48pt}{7.15pt}{E}};
\node at (206.44,124.70) {\normalfont\fontsize{11}{11}\bfseries\itshape\selectfont \resizebox{6.40pt}{7.27pt}{D}};
\node at (71.92,181.02) {\normalfont\fontsize{11}{11}\bfseries\itshape\selectfont \resizebox{6.53pt}{6.90pt}{N}};
\node at (127.15,203.40) {\normalfont\fontsize{11}{11}\bfseries\itshape\selectfont \resizebox{7.31pt}{7.51pt}{C}};
\end{tikzpicture}
\end{minipage}\hfill\begin{minipage}[t]{0.589\FullSourceWidth}\vspace{0pt}\setlength{\GutterOffset}{0.411\FullSourceWidth}
\par\vspace{3.30mm}
\hypertarget{AB01-PDF0308-P04}{}
\DLine{AB01-PDF0308-body-L024}{}{}{\Indent \textit{Aequationem centri} supputare.}
\hypertarget{AB01-PDF0308-P05}{}
\DLine{AB01-PDF0308-body-L025}{}{}{\Indent Constructiones geometricae p. 46–47 ex}
\DLine{AB01-PDF0308-body-L026}{}{}{\textit{Almag.} III, 5 (ed. Halma t. I, p. 190–199)}
\par\vspace{0.09mm}
\DLine{AB01-PDF0308-body-L027}{}{}{sumptae sunt; deest contra apud Ptolemaeum}
\par\vspace{0.14mm}
\DLine{AB01-PDF0308-body-L028}{}{}{formula \Name{al–Battānī}:}
\par\vspace{1.85mm}
\DLine{AB01-PDF0308-body-L029}{}{}{\mbox{}\rlap{\hspace*{1.64mm}sin aeq. centri $= \dfrac{R\,\dfrac{e \sin m}{R}}{\sqrt{\Bigl(\dfrac{e \sin m}{R}\Bigr)^{2} + \Bigl(\dfrac{e \cos m}{R} + 60^{p}\Bigr)^{2}}}$\,,}}
\par\vspace{0.22mm}
\DLine{AB01-PDF0308-body-L030}{}{}{ubi \textit{e} est excentricitas vel, in hypothesi epi-}
\DLine{AB01-PDF0308-body-L031}{}{}{cycli, radius epicycli (1); \textit{m} arcus anomaliae,}
\DLine{AB01-PDF0308-body-L032}{}{}{scilicet arcus excentrici vel epicycli quem sidus}
\par\vspace{0.11mm}
\DLine{AB01-PDF0308-body-L033}{}{}{ab apogeo usque ad datum temporis momen-}
\DLine{AB01-PDF0308-body-L034}{}{}{tum percurrerit. Facillime regula deducitur e}
\end{minipage}\par
\par\vspace{0.76mm}
\par\vspace{6pt}\nointerlineskip\hbox to\textwidth{\hss\vrule width 115.43pt height.25pt depth0pt\hss}\prevdepth=-1000pt \vspace{5pt}
\noindent\begin{minipage}[t]{.48\textwidth}\vspace{0pt}\fontsize{9}{11.5}\selectfont\setlength{\GutterOffset}{0pt}\setlength{\FullSourceWidth}{\linewidth}
\par\fontsize{9}{11.80}\selectfont
\par\vspace{0.04mm}
\DLine{AB01-PDF0308-notes_left-L001}{}{}{supputassent, longitudinem apogei 79° 30′ ponebant.}
\DLine{AB01-PDF0308-notes_left-L002}{}{}{Est quantitas recepta in eorum \textit{Tabula Probata} et}
\DLine{AB01-PDF0308-notes_left-L003}{}{}{ab \Name{Ḥabash}, f. 28,r. — Perperam \Name{al-Mas‘ūdī},}
\end{minipage}
\hfill\vrule width.3pt\hfill\begin{minipage}[t]{.48\textwidth}\vspace{0pt}\fontsize{9}{11.5}\selectfont\setlength{\GutterOffset}{0pt}\setlength{\FullSourceWidth}{\linewidth}
\par\fontsize{9}{11.80}\selectfont
\par\vspace{-0.09mm}
\DLine{AB01-PDF0308-notes_right-L001}{}{}{\textit{at-Tanbīh} p. 222, hanc observationem anno 217 heg.}
\par\vspace{-0.13mm}
\DLine{AB01-PDF0308-notes_right-L002}{}{}{fuisse scribit.}
\hypertarget{AB01-PDF0308-N01}{}
\DLine{AB01-PDF0308-notes_right-L003}{}{}{\Indent (1) In partibus radii excentrici vel orbitae ($=$60\textsuperscript{\textit{p}}).}
\end{minipage}
\end{document}
